\documentclass[11pt]{ctexart}
\usepackage[margin=1in]{geometry}
\usepackage{amsmath,amssymb,booktabs}
\usepackage{enumitem}
\usepackage[hidelinks]{hyperref}

\title{论文总结：{{PAPER_TITLE}}}
\author{Generated by literature-summary-agent}
\date{{{REPORT_DATE}}}

\begin{document}
\maketitle

\section{基本信息}
\begin{itemize}[leftmargin=1.5em]
  \item 标题：{{PAPER_TITLE}}
  \item 作者：{{AUTHORS}}
  \item 年份/发表：{{YEAR_VENUE}}
  \item Zotero 条目 Key：{{ITEM_KEY}}
  \item PDF 路径：\texttt{{{PDF_PATH}}}
\end{itemize}

\section{第一层：问题逻辑}
\subsection{问题定义}
{{PROBLEM}}

\subsection{现有方法缺口}
{{GAP}}

\subsection{从问题到创新的推理链}
{{REASONING}}

\subsection{创新总览}
{{INNOVATION_SUMMARY}}

\section{第二层：创新结构}
\subsection{创新点一：{{INNOVATION1_NAME}}}
\paragraph{原理}
{{INNOVATION1_PRINCIPLE}}
\paragraph{实现}
{{INNOVATION1_IMPLEMENTATION}}
\paragraph{效果}
{{INNOVATION1_EFFECT}}

\subsection{创新点二：{{INNOVATION2_NAME}}}
\paragraph{原理}
{{INNOVATION2_PRINCIPLE}}
\paragraph{实现}
{{INNOVATION2_IMPLEMENTATION}}
\paragraph{效果}
{{INNOVATION2_EFFECT}}

\section{必要公式}
\subsection{公式一}
\begin{equation}
{{EQUATION1}}
\end{equation}
{{EQUATION1_EXPLAIN}}

\subsection{公式二}
\begin{equation}
{{EQUATION2}}
\end{equation}
{{EQUATION2_EXPLAIN}}

\section{第三层：证据链验证}
\subsection{整体对比}
{{OVERALL_COMPARISON}}

\subsection{消融证据}
{{ABLATION_EVIDENCE}}

\subsection{假设验证}
{{HYPOTHESIS_VALIDATION}}

\subsection{因果主张检查}
{{CAUSAL_CHECK}}

\section{一句话结论}
{{ONE_SENTENCE_CONCLUSION}}

\end{document}
